% !PRINTED 416
\weberpage{416}
%<<CONT
consists, which is satisfied for all values of the $x$ that satisfy
the relations (1), and in which $z$, always with regard to the
relations (1), remains unchanged when the $x$ are subjected to the
permutations of the group $\mathfrak{A}$. The function $u$ is thereby
to pass into $\nu$ distinct functions
\weq{3}{u, u_1, \ldots, u_{\nu-1}}
so that the quantities (3) are the roots of the equation (2), and,
when $\nu > 1$, (2) is a resolvent of the equation for the $x$.

In this generality we have as yet no point of attack for the solution
of the question. We therefore add a further restrictive assumption,
namely: \gesperrt{It is to be possible to set, for the $x_h$, rational
functions of a variable $t$}
\weq{4}{x_h = \varphi_h(t)}
\gesperrt{so that the relations (1) are identically satisfied, and that at the
same time $z$ does not become independent of
$t$}\footnote{If one regards the $x$ and the $u$ as the coordinates
each of a point $X$ and $U$ in a space of $n$ and of $\nu$
dimensions, then the relations (1) determine for $x$ a manifold of
fewer than $n$ dimensions. To each point $X$ there corresponds a
point $U$, and to the totality of the $X$, by (2), a manifold of the
$U$ of one dimension (a curve). Our assumption is that this is a
so-called rational or unicursal curve, or (in the language of the
theory of functions) a curve of genus zero.}.

To give an example here at once, one with which we shall later occupy
ourselves in detail, we mention the case of five variables
$x_0, x_1, x_2, x_3, x_4$ which satisfy the conditions
\weq{5}{\Sigma x = 0, \quad \Sigma x^2 = 0}
and which are therefore the roots of a \gesperrt{principal equation} of degree 5
(Bd. I, \S.\,54). That the assumption we have made holds in this
case is shown by the representation in Bd. I, \S.\,74, (4):
\[
y = t_3 F_0 + t_2(\alpha_2 F_1 + \alpha_1 F_2 + \alpha_0 F_3),
\]
which, when the determinations given there are made, put for the $x$,
satisfies the conditions (5) for all values of the ratio
$t = t_2 : t_3$.
